\chapter{模60计数器}

\section{应用统一框架}

\begin{enumerate}
  \item \textbf{模值N}：$N = 60$
  \item \textbf{寄存器位宽k}：$k = \lceil \log_2 60 \rceil = 6$（$2^5=32<60<2^6=64$）
  \item \textbf{状态转移}：count $\leftarrow$ count + 1
  \item \textbf{回零条件}：count = 59(111011)时，强制归零
\end{enumerate}

\section{模60的实际意义}

模60计数器是数字钟表的核心：
\begin{itemize}
  \item 秒：0-59（模60）
  \item 分：0-59（模60）
\end{itemize}

一个完整的数字钟 = 模60秒 + 模60分 + 模24（或模12）时。

\section{VHDL实现}

\begin{lstlisting}[caption={模60二进制计数器}]
library ieee;
use ieee.std_logic_1164.all;
use ieee.std_logic_unsigned.all;

entity counter_mod60 is
  port (
    clk   : in  std_logic;
    rst_n : in  std_logic;
    count : out std_logic_vector(5 downto 0);
    carry : out std_logic
  );
end counter_mod60;

architecture behavioral of counter_mod60 is
  signal cnt : std_logic_vector(5 downto 0);
begin
  process(clk, rst_n)
  begin
    if rst_n = '0' then
      cnt <= (others => '0');
    elsif rising_edge(clk) then
      if cnt = "111011" then    -- = 59
        cnt <= (others => '0');
      else
        cnt <= cnt + 1;
      end if;
    end if;
  end process;

  count <= cnt;
  carry <= '1' when cnt = "111011" else '0';
end behavioral;
\end{lstlisting}

\section{三种模值计数器的统一对比}

\begin{center}
\begin{tabular}{|l|c|c|c|}
\hline
\textbf{参数} & \textbf{模12} & \textbf{模24} & \textbf{模60} \\
\hline
模值N & 12 & 24 & 60 \\
位宽 & 4位 & 5位 & 6位 \\
$2^k$ & 16 & 32 & 64 \\
冗余状态数 & 4 & 8 & 4 \\
清零条件 & cnt=11 & cnt=23 & cnt=59 \\
清零二进制 & 1011 & 10111 & 111011 \\
\hline
VHDL差异 & \multicolumn{3}{c|}{\textbf{仅修改：位宽 + 清零比较值}} \\
\hline
\end{tabular}
\end{center}

\begin{keypoint}[所有二进制计数器的统一模板]
\textbf{修改任意模N计数器的三个步骤：}

1. 计算位宽：$k = \lceil \log_2 N \rceil$
2. 修改信号位宽为k
3. 修改清零条件为：cnt = N-1

\textbf{其余逻辑完全不变。}
\end{keypoint}
